% ---------------------------------------------------------------------------
% Banca 1 — deck Beamer (tema Warsaw)
%
% Artefato DERIVADO de docs/07_apresentacoes/banca1/02_conteudo_slides.md.
% Divergencia entre este deck e o documento de conteudo e erro do deck.
%
% Compilar SEMPRE a partir da raiz do repositorio:
%     bash scripts/apresentacao/build_deck_beamer.sh
% ---------------------------------------------------------------------------
\documentclass[aspectratio=169,11pt]{beamer}

\usepackage[T1]{fontenc}
\usepackage[utf8]{inputenc}
\usepackage[brazil]{babel}
\usepackage{lmodern}
\usepackage{amsmath}
\usepackage{amssymb}
\usepackage{booktabs}
\usepackage{graphicx}
\usepackage{array}
\usepackage{tabularx}
\usepackage{makecell}
\usepackage{ragged2e}

% --- figuras: caminhos relativos a raiz do repositorio ----------------------
\graphicspath{{output/apresentacao_banca1/}{docs/02_teoria/figuras/}}

% --- tema -------------------------------------------------------------------
\usetheme{Warsaw}
\setbeamertemplate{navigation symbols}{}

\definecolor{verdeEscuro}{HTML}{16301F}
\definecolor{verdeMeio}{HTML}{2A5539}
\definecolor{verdeMedio}{HTML}{8AAA92}
\definecolor{verdeClaro}{HTML}{DDE7DE}
\definecolor{verdePalido}{HTML}{F0F5F0}
\definecolor{cinzaTexto}{HTML}{2E2E2E}
\definecolor{cinzaFonte}{HTML}{6B6B6B}

\setbeamercolor{normal text}{fg=cinzaTexto,bg=white}
\setbeamercolor{structure}{fg=verdeMeio}
\setbeamercolor{alerted text}{fg=verdeMeio}

% frametitle do Warsaw e um degrade entre estas duas cores
\setbeamercolor{subsection in head/foot}{bg=verdeMeio,fg=white}
\setbeamercolor{section in head/foot}{bg=verdeEscuro,fg=white}

\setbeamercolor{title}{bg=verdeEscuro,fg=white}
\setbeamercolor{subtitle}{parent=title,fg=verdeClaro}
\setbeamercolor{author}{fg=cinzaFonte}
\setbeamercolor{block title}{bg=verdeMeio,fg=white}
\setbeamercolor{block body}{bg=verdePalido,fg=cinzaTexto}
\setbeamercolor{itemize item}{fg=verdeMeio}
\setbeamercolor{itemize subitem}{fg=verdeMedio}
\setbeamercolor{enumerate item}{fg=verdeMeio}
\setbeamercolor{rastreio}{bg=verdeClaro,fg=verdeEscuro}
\setbeamercolor{rodape}{bg=verdeEscuro,fg=white}

\setbeamerfont{frametitle}{size=\large,series=\bfseries}
\setbeamerfont{title}{size=\large,series=\bfseries}
\setbeamerfont{subtitle}{size=\small}
\setbeamerfont{block title}{size=\normalsize,series=\bfseries}
\setbeamerfont{itemize/enumerate body}{size=\normalsize}
\setbeamerfont{itemize/enumerate subbody}{size=\small}

\setbeamertemplate{itemize item}{\raisebox{0.15ex}{\scriptsize$\blacktriangleright$}}
\setbeamertemplate{itemize subitem}{\raisebox{0.15ex}{\scriptsize$\bullet$}}

\setlength{\leftmargini}{1.4em}
\setlength{\leftmarginii}{1.2em}

% --- faixa fina de rastreio de secao (fora do titulo) -----------------------
\def\rastreiotexto{}
\newcommand{\rastreio}[1]{\gdef\rastreiotexto{#1}}

\setbeamertemplate{headline}{%
  \ifx\rastreiotexto\empty
    \vskip0pt%
  \else
    \begin{beamercolorbox}[wd=\paperwidth,ht=2.6ex,dp=1.4ex,%
      leftskip=1.1em,rightskip=1.1em]{rastreio}%
      \scriptsize\rastreiotexto\hfill%
    \end{beamercolorbox}%
  \fi
}

\setbeamertemplate{footline}{%
  \begin{beamercolorbox}[wd=\paperwidth,ht=2.6ex,dp=1.4ex,%
    leftskip=1.1em,rightskip=1.1em]{rodape}%
    \scriptsize Banca 1 \textperiodcentered\ Desvantagens territoriais na escolha
    locacional de médicos especialistas%
    \hfill\insertframenumber\,/\,\inserttotalframenumber%
  \end{beamercolorbox}%
}

% --- utilitarios de composicao ---------------------------------------------
% Unico formato de fonte no deck: filete fino, rotulo em negrito, alinhado a
% esquerda, corpo reduzido. Vale tambem para os frames com figura.
\newcommand{\fonte}[1]{%
  \vfill
  {\color{verdeMedio}\rule{\textwidth}{0.4pt}}\par
  \vspace{0.25em}
  {\color{cinzaFonte}\scriptsize\setlength{\baselineskip}{1.05\baselineskip}#1\par}%
}

\newcommand{\lead}[1]{{\color{verdeEscuro}\bfseries #1}}
\newcommand{\num}[1]{{\color{verdeMeio}\bfseries #1}}

\newcolumntype{L}{>{\raggedright\arraybackslash}X}
\newcolumntype{Z}{>{\centering\arraybackslash}X}

\renewcommand{\arraystretch}{1.28}

\setbeamertemplate{caption}{\insertcaption}

% --- metadados --------------------------------------------------------------
\title{Desvantagens territoriais na escolha locacional de médicos especialistas}
\subtitle{Um modelo microeconômico para o preenchimento de vagas do Programa
Mais Médicos Especialistas}
\author{Autoria \textperiodcentered\ instituição \textperiodcentered\ data da banca}
\date{}

\begin{document}

% ===========================================================================
% 1 — Capa
% ===========================================================================
\rastreio{}
\begin{frame}[plain]
  \vfill
  \begin{center}
    \begin{beamercolorbox}[wd=0.82\textwidth,sep=0.9em,center,rounded=true,%
      shadow=true]{title}
      \usebeamerfont{title}\inserttitle\par
      \vspace{0.45em}
      {\usebeamerfont{subtitle}\usebeamercolor[fg]{subtitle}\insertsubtitle\par}
    \end{beamercolorbox}

    \vspace{1.1em}
    {\color{cinzaFonte}\small\insertauthor\par}
  \end{center}
  \vfill
\end{frame}

% ===========================================================================
% 2 — Sumário
% ===========================================================================
\rastreio{}
\begin{frame}{Sumário}
  \vfill
  \begin{columns}[T]
    \begin{column}{0.48\textwidth}
      \large
      \begin{enumerate}
        \setlength{\itemsep}{1.1em}
        \item Motivação
        \item Pergunta
        \item Literatura teórica
      \end{enumerate}
    \end{column}
    \begin{column}{0.48\textwidth}
      \large
      \begin{enumerate}
        \setlength{\itemsep}{1.1em}
        \setcounter{enumi}{3}
        \item Modelo microeconômico
        \item Hipóteses
        \item Viabilidade empírica
      \end{enumerate}
    \end{column}
  \end{columns}
  \vfill
\end{frame}

% ===========================================================================
% 3 — Especialistas não faltam; faltam no interior      (2 frames)
% ===========================================================================
\rastreio{1. Motivação \textperiodcentered\ Problema \textperiodcentered\ 1 de 3}

\begin{frame}{Especialistas não faltam; faltam no interior}
  \vspace{0.2em}
  O Brasil tinha, em 2024, \num{353 mil médicos especialistas} --- 59\% dos
  597 mil médicos do país. \lead{O problema não é o número. É onde eles estão.}

  \vspace{0.9em}
  \begin{columns}[T]
    \begin{column}{0.46\textwidth}
      \begin{beamercolorbox}[wd=\textwidth,sep=0.7em,rounded=true]{block body}
        {\Large\bfseries\color{verdeEscuro}55\%} dos especialistas estão no
        \textbf{Sudeste};\\[0.25em]
        {\Large\bfseries\color{verdeEscuro}6\%} no \textbf{Norte}
      \end{beamercolorbox}
    \end{column}
    \begin{column}{0.50\textwidth}
      \begin{beamercolorbox}[wd=\textwidth,sep=0.7em,rounded=true]{block body}
        {\Large\bfseries\color{verdeEscuro}453} especialistas por 100 mil
        habitantes no \textbf{Distrito Federal};\\[0.25em]
        {\Large\bfseries\color{verdeEscuro}68} no \textbf{Maranhão} e
        {\Large\bfseries\color{verdeEscuro}70} no \textbf{Pará}
      \end{beamercolorbox}
    \end{column}
  \end{columns}

  \vspace{0.9em}
  {\small\itshape ``apenas 10\% dos especialistas atendem no SUS. Além disso,
  há concentração desses profissionais nas capitais e regiões mais ricas do
  país''}

  \fonte{\textbf{Fontes:} Scheffer et al., \textit{Demografia Médica no Brasil
  2025} (FMUSP/AMB), dados de dezembro de 2024; citação de Senado Notícias,
  25/09/2025, referindo dados do Ministério da Saúde.}
\end{frame}

\begin{frame}{Especialistas não faltam; faltam no interior}
  \vspace{0.3em}
  \lead{Quem sente a falta é o paciente do SUS fora dos grandes centros.}

  \vspace{0.7em}
  Em 2025 o Ministério da Saúde reconheceu \num{situação de urgência em saúde
  pública} em todo o país, por \textbf{24 meses}, em razão do tempo de espera
  por consultas, exames e cirurgias na atenção especializada --- e lançou o
  programa \textbf{Agora Tem Especialistas}, do qual o \textbf{PMM-E} é o braço
  de provimento.

  \vspace{1.0em}
  \begin{columns}[T]
    \begin{column}{0.48\textwidth}
      \begin{beamercolorbox}[wd=\textwidth,sep=0.7em,rounded=true]{block body}
        \small\itshape ``Ministério da Saúde decreta situação de urgência em
        saúde pública pelos próximos dois anos''
      \end{beamercolorbox}
    \end{column}
    \begin{column}{0.48\textwidth}
      \begin{beamercolorbox}[wd=\textwidth,sep=0.7em,rounded=true]{block body}
        \small\itshape ``Aprovada no Senado, MP deve reduzir fila para
        atendimento por especialista no SUS''
      \end{beamercolorbox}
    \end{column}
  \end{columns}

  \fonte{\textbf{Fontes:} Portaria GM/MS nº 7.061, de 6 de junho de 2025;
  manchetes do Correio do Povo, 07/05/2025, e de Senado Notícias, 25/09/2025.}
\end{frame}

% ===========================================================================
% 4 — Onde a bolsa é maior, já havia menos especialistas   (2 frames)
% ===========================================================================
\rastreio{1. Motivação \textperiodcentered\ Problema \textperiodcentered\ 2 de 3}

\begin{frame}{Onde a bolsa é maior, já havia menos especialistas}
  \vspace{0.2em}
  {\small O retrato nacional se repete dentro do programa: nos \textbf{295
  municípios} com vaga no ciclo 1, um mês antes da oferta, \lead{menos da
  metade} de especialistas por habitante onde a bolsa seria maior.}

  \vspace{0.3em}
  \begin{center}
    \includegraphics[width=\textwidth,height=0.58\textheight,keepaspectratio]%
      {oferta_pre_por_faixa}
  \end{center}

  \fonte{\textbf{Fontes:} CNES 06/2025 --- profissionais nos CBOs dos 10 cursos
  com correspondência unívoca curso--CBO, nos 295 municípios com vaga no ciclo
  1; população do Censo 2022 (IBGE); faixa pela categoria de IVS do Ipea.}
\end{frame}

\begin{frame}{Onde a bolsa é maior, já havia menos especialistas}
  \vspace{0.2em}
  {\small\lead{E quem vai, vai sozinho.} Na Faixa 1, o especialista encontraria
  menos colegas da sua especialidade --- e é lá que é maior a chance de ser o
  único, ou de ter um só colega.}

  \vspace{0.3em}
  \begin{center}
    \includegraphics[width=\textwidth,height=0.58\textheight,keepaspectratio]%
      {retaguarda_por_faixa}
  \end{center}

  \fonte{\textbf{Fontes:} CNES 06/2025, nos 295 municípios com vaga no ciclo 1;
  faixa pela categoria de IVS do Ipea. A Faixa 1 tem 19 pares
  município--especialidade: a porcentagem oscila com poucos casos, a mediana
  não.}
\end{frame}

% ===========================================================================
% 5 — O que o médico vê ao decidir                        (2 frames)
% ===========================================================================
\rastreio{1. Motivação \textperiodcentered\ Problema \textperiodcentered\ 3 de 3}

\begin{frame}{O que o médico vê ao decidir}
  \vspace{0.2em}
  {\small Para o médico, ``município vulnerável'' não é um índice. É um
  conjunto de \textbf{desvantagens concretas}. Quatro aparecem de forma
  consistente na literatura --- \lead{duas delas nós conseguimos medir}.}

  \vspace{0.6em}
  {\footnotesize\renewcommand{\arraystretch}{1.18}
  \begin{tabularx}{\textwidth}{@{}>{\bfseries\raggedright\arraybackslash}p{0.20\textwidth} L >{\raggedright\arraybackslash}p{0.26\textwidth}@{}}
    \toprule
    \normalfont\bfseries Desvantagem & \bfseries O que significa para o médico
      & \bfseries Medimos? \\
    \midrule
    Retaguarda profissional & poucos ou nenhum colega da mesma especialidade;
      sem segunda opinião, sem escala de plantão, sem a quem encaminhar o caso
      difícil & \textbf{sim} --- slide anterior \\
    Infraestrutura & equipamento, insumos e equipe de apoio escassos; o
      atendimento cansa mais e resolve menos & \textbf{em parte} --- leitos e
      equipamentos do CNES \\
    Distância da família & o custo de viver longe de onde a família está, e
      onde o médico nasceu ou se formou & \textbf{não} \\
    Mercado privado ausente & sem consultório ou plano de saúde local, a
      remuneração se reduz ao que o programa paga & \textbf{não} \\
    \bottomrule
  \end{tabularx}}
\end{frame}

\begin{frame}{O que o médico vê ao decidir}
  \vspace{0.2em}
  \lead{O que a literatura diz sobre o peso de cada uma}

  \vspace{0.7em}
  {\small
  \begin{itemize}
    \setlength{\itemsep}{0.95em}
    \item \textbf{Estados Unidos}, início do século XX: a escolha do lugar já
      dependia de \textit{``preferences over rural or urban living, or other
      location-specific attributes, such as proximity to family''}.
    \item \textbf{Brasil}, 50 mil generalistas formados de 2001 a 2013:
      \num{a proximidade do lugar de nascimento ou de formação é o principal
      fator} da escolha locacional; salário e infraestrutura importam, mas em
      escala menor.
    \item \textbf{Austrália}, 3.727 clínicos perguntados sobre o que os faria
      mudar para o interior: \num{65\% não mudariam por pacote nenhum}; para os
      demais, o valor exigido depende do tamanho da cidade e --- nas palavras
      dos autores --- \textit{``not only of the area but also of the
      characteristics of the job''}.
  \end{itemize}}

  \fonte{\textbf{Fontes:} Moehling, Niemesh, Thomasson \& Treber (2020),
  \textit{Cliometrica}, p. 184; Costa, Nunes \& Sanches (2024), \textit{Review
  of Economics and Statistics}; Scott et al. (2013), \textit{Social Science \&
  Medicine}.}
\end{frame}

% ===========================================================================
% 6 — O que é o PMM-E                                     (4 frames)
% ===========================================================================
\rastreio{1. Motivação \textperiodcentered\ Política \textperiodcentered\ 1 de 2}

\begin{frame}{O que é o PMM-E}
  \vspace{0.4em}
  \begin{block}{Lei}
    A \textbf{Lei nº 15.233/2025} criou o \textbf{Projeto Mais Médicos
    Especialistas} dentro do Programa Mais Médicos, \textit{``destinado ao
    provimento de profissionais com vistas à redução no tempo de espera de
    atendimento ao usuário do SUS, nas regiões prioritárias''}, como braço
    formativo do programa \textbf{Agora Tem Especialistas}.
  \end{block}

  \vspace{0.6em}
  \begin{block}{Quem}
    Médicos com diploma brasileiro ou revalidado e \textbf{registro de
    especialista (RQE)} na área da vaga. Não é concurso nem emprego: o médico
    recebe \num{bolsa-formação} mensal do Ministério da Saúde, \textbf{sem
    vínculo}.
  \end{block}

  \fonte{\textbf{Fontes:} Lei nº 15.233/2025, art. 22-D; Portaria GM/MS
  nº 7.177/2025; Edital SGTES/MS nº 3/2025 (DOU 24/07/2025), itens 1 e 3.}
\end{frame}

\begin{frame}{O que é o PMM-E}
  \vspace{0.3em}
  \lead{O quê.} Um \textbf{aprimoramento em serviço} de \num{12 meses}, com
  \num{20 horas semanais} em estabelecimento do SUS, supervisão e mentoria de
  uma instituição formadora, e imersões em serviços de referência.

  \vspace{0.7em}
  {\small São \textbf{16 cursos}:}
  \vspace{0.3em}

  \begin{columns}[T]
    \begin{column}{0.47\textwidth}
      \begin{beamercolorbox}[wd=\textwidth,sep=0.7em,rounded=true]{block body}
        \small\textbf{6 cirúrgicos}\\[0.3em]
        anestesiologia; cirurgia geral, oncológica, colorretal, digestiva e
        ginecológica
      \end{beamercolorbox}
    \end{column}
    \begin{column}{0.47\textwidth}
      \begin{beamercolorbox}[wd=\textwidth,sep=0.7em,rounded=true]{block body}
        \small\textbf{10 ambulatoriais}\\[0.3em]
        endoscopia e colonoscopia, oncologia clínica, radioterapia,
        ecocardiografia, ultrassonografia mamária, colposcopia,
        videolaringoscopia, anatomia patológica
      \end{beamercolorbox}
    \end{column}
  \end{columns}

  \vspace{0.8em}
  \centering\lead{O foco é o câncer e o diagnóstico que o SUS mais espera.}

  \fonte{\textbf{Fontes:} Edital SGTES/MS nº 3/2025 (DOU 24/07/2025),
  itens 1, 3, 4, 5, 10 e 11.}
\end{frame}

\begin{frame}{O que é o PMM-E}
  \vspace{0.4em}
  \lead{Como a vaga chega ao médico}

  \vspace{0.9em}
  {\small
  \begin{enumerate}
    \setlength{\itemsep}{0.8em}
    \item O \textbf{estado ou município} indica o serviço e a especialidade.
    \item A \textbf{comissão bipartite} prioriza.
    \item O \textbf{Ministério} analisa a capacidade instalada e publica o
      quadro de vagas com município, estabelecimento, curso e \textbf{faixa de
      bolsa}.
    \item O \textbf{médico} escolhe \textbf{até dois locais}, em ordem de
      preferência, e é classificado por titulação e tempo de formação.
  \end{enumerate}}

  \vspace{0.9em}
  \centering
  \begin{beamercolorbox}[wd=0.86\textwidth,sep=0.7em,rounded=true]{block body}
    \small O serviço \textbf{não pode substituir} profissional já contratado
    por um bolsista.
  \end{beamercolorbox}

  \fonte{\textbf{Fontes:} Portaria GM/MS nº 7.177/2025; Edital SGTES/MS
  nº 3/2025 (DOU 24/07/2025), itens 4 e 5.}
\end{frame}

\begin{frame}{O que é o PMM-E}
  \vspace{0.4em}
  \begin{columns}[T]
    \begin{column}{0.31\textwidth}
      {\footnotesize\lead{Onde, no primeiro ciclo (julho de 2025).}

      \vspace{0.5em}
      \num{1.295 vagas} estabelecimento--curso em \num{460 estabelecimentos}
      e \num{368 municípios}, em todas as UFs --- 678 para preenchimento
      imediato e 1.145 em cadastro de reserva.

      \vspace{0.5em}
      O \textbf{Nordeste} concentra 39\% das vagas; \textbf{Minas Gerais} é o
      estado com mais vagas (252).

      \vspace{0.5em}
      Dois terços dos municípios têm menos de 100 mil habitantes;
      \textbf{18 são capitais}.}
    \end{column}
    \begin{column}{0.65\textwidth}
      \includegraphics[width=\textwidth,height=0.66\textheight,keepaspectratio]%
        {vagas_ciclo1_por_regiao}
    \end{column}
  \end{columns}

  \fonte{\textbf{Fontes:} quadro de vagas do ciclo 1, chamada 1;
  Censo 2022 (IBGE).}
\end{frame}

% ===========================================================================
% 7 — A bolsa remunera o lugar                            (3 frames)
% ===========================================================================
\rastreio{1. Motivação \textperiodcentered\ Política \textperiodcentered\ 2 de 2}

\begin{frame}{A bolsa remunera o lugar}
  \vspace{0.2em}
  {\small O valor da bolsa \textbf{não} depende da especialidade, da carga nem
  do que o médico produz. Depende de \lead{onde fica o município}, em três
  passos:}

  \vspace{0.7em}
  {\small
  \begin{enumerate}
    \setlength{\itemsep}{0.7em}
    \item \textbf{O índice.} O \textbf{IVS 2010 do Ipea} resume, em um número
      de 0 a 1, dezesseis indicadores do Censo 2010 em três dimensões:
      \textit{infraestrutura urbana} (saneamento, lixo, tempo de deslocamento),
      \textit{capital humano} (mortalidade infantil, analfabetismo, crianças
      fora da escola) e \textit{renda e trabalho} (pobreza, desemprego,
      informalidade).
    \item \textbf{A categoria.} O Ipea corta o índice em cinco categorias:
      muito baixa (até 0,200), baixa (0,201--0,300), média (0,301--0,400),
      alta (0,401--0,500) e muito alta (acima de 0,500).
    \item \textbf{O valor.} O edital de 2025 agrupou as cinco categorias em
      \textbf{três faixas}.
  \end{enumerate}}

  \fonte{\textbf{Fontes:} Ipea, \textit{Atlas da Vulnerabilidade Social nos
  Municípios Brasileiros} (2015); Edital SGTES/MS nº 3/2025, item 11.1.3.}
\end{frame}

\begin{frame}{A bolsa remunera o lugar}
  \vspace{0.5em}
  \begin{columns}[c]
    \begin{column}{0.44\textwidth}
      {\footnotesize
      \begin{tabular}{@{}l c r@{}}
        \toprule
        \makecell[l]{\bfseries Categoria\\\bfseries de IVS} &
        \bfseries Faixa &
        \makecell[r]{\bfseries Bolsa\\\bfseries mensal} \\
        \midrule
        muito alta & 1 & R\$ 20.000 \\
        alta       & 2 & R\$ 15.000 \\
        \makecell[l]{média, baixa\\ou muito baixa} & 3 & R\$ 10.000 \\
        \bottomrule
      \end{tabular}}
    \end{column}
    \begin{column}{0.53\textwidth}
      \includegraphics[width=\textwidth,height=0.50\textheight,keepaspectratio]%
        {bolsa_por_faixa}
    \end{column}
  \end{columns}

  \vspace{1.0em}
  {\small No ciclo 1, \num{102} municípios foram publicados na Faixa 1,
  \num{107} na Faixa 2 e \num{159} na Faixa 3.}

  \fonte{\textbf{Fontes:} Edital SGTES/MS nº 3/2025, item 11.1.3, e
  retificação; quadro de vagas do ciclo 1.}
\end{frame}

\begin{frame}{A bolsa remunera o lugar}
  \vspace{0.6em}
  \lead{Dois cuidados}

  \vspace{1.0em}
  \begin{block}{A categoria que vale é a publicada na vaga}
    Em \num{177 dos 368 municípios} ela \textbf{não coincide} com a que se
    obtém aplicando os cortes do Ipea ao IVS local.
  \end{block}

  \vspace{0.8em}
  \begin{block}{A grade mudou em 2026}
    A categoria \textit{alta} passou à \textbf{Faixa 1}.
  \end{block}

  \fonte{\textbf{Fontes:} Edital SGTES/MS nº 3/2025, item 11.1.3, e
  retificação; Chamamento SGTES/MS nº 1/2026; Ipea, \textit{Atlas da
  Vulnerabilidade Social nos Municípios Brasileiros} (2015); quadro de vagas do
  ciclo 1.}
\end{frame}

% ===========================================================================
% 8 — Pagar mais funciona: a evidência a favor            (2 frames)
% ===========================================================================
\rastreio{1. Motivação \textperiodcentered\ Efeito incerto \textperiodcentered\ 1 de 3}

\begin{frame}{Pagar mais funciona: a evidência a favor}
  \vspace{0.3em}
  {\small O programa aposta que dinheiro compensa lugar ruim.
  \lead{A aposta tem precedente.}}

  \vspace{0.9em}
  \begin{block}{No México, o salário foi sorteado}
    \small
    Um concurso público real distribuiu \textbf{106 postos} em municípios
    pobres e anunciou, \textbf{ao acaso}, dois salários.

    \vspace{0.5em}
    \begin{itemize}
      \setlength{\itemsep}{0.45em}
      \item Onde o salário era \textbf{33\% maior}, a aceitação da vaga subiu
        \num{15 pontos percentuais}.
      \item O efeito foi maior onde o lugar era pior: nos postos a mais de
        \textbf{200 km} da cidade natal do candidato, a aceitação foi de
        \num{25\% para cerca de 80\%}.
      \item O aumento \textbf{anulou} a rejeição aos municípios de menor
        desenvolvimento humano, \textbf{sem} atrair candidatos menos
        qualificados ou menos motivados.
    \end{itemize}
  \end{block}

  \fonte{\textbf{Fonte:} Dal Bó, Finan \& Rossi (2013), \textit{Quarterly
  Journal of Economics}.}
\end{frame}

\begin{frame}{Pagar mais funciona: a evidência a favor}
  \vspace{0.3em}
  \lead{O degrau do PMM-E tem o tamanho que a literatura pede.}

  \vspace{0.9em}
  {\small
  Os estudos que medem quanto um médico exige para ir a um posto pior chegam a
  valores entre \num{37\% e 64\%} da renda anual para cidades pequenas
  (Austrália), e a oferta de médicos no interior brasileiro responde a salário
  com elasticidade da ordem de \num{0,7}.}

  \vspace{0.9em}
  \centering
  \begin{beamercolorbox}[wd=0.80\textwidth,sep=0.8em,rounded=true,shadow=true]{block title}
    \centering O degrau do programa --- R\$ 5 mil sobre R\$ 10 mil, ou
    \textbf{+50\%} --- cai dentro dessa faixa.
  \end{beamercolorbox}

  \vspace{1.0em}
  \raggedright
  {\small\textbf{Ressalva:} os percentuais da literatura são sobre a
  \textbf{renda total} do médico; a bolsa do PMM-E remunera \textbf{20 horas
  semanais}.}

  \fonte{\textbf{Fontes:} Scott et al. (2013), \textit{Social Science \&
  Medicine}; Costa, Nunes \& Sanches (2024), \textit{Review of Economics and
  Statistics}.}
\end{frame}

% ===========================================================================
% 9 — Mas é caro, e não segura: a evidência contra        (2 frames)
% ===========================================================================
\rastreio{1. Motivação \textperiodcentered\ Efeito incerto \textperiodcentered\ 2 de 3}

\begin{frame}{Mas é caro, e não segura: a evidência contra}
  \vspace{0.3em}
  \begin{block}{Muitos não vão por preço nenhum}
    \small
    Dos \textbf{3.727 clínicos australianos}, \num{65\%} escolheram ficar onde
    estavam em todos os cenários oferecidos. Para o pior posto, quem mudaria
    pedia \num{130\%} da renda anual.
  \end{block}

  \vspace{0.7em}
  \begin{block}{No Brasil, salário compra pouco e custa muito}
    \small
    Um modelo calibrado com todos os generalistas formados entre 2001 e 2013
    estima que aumentar em \textbf{50\%} o salário público no interior do Norte
    e do Nordeste corrigiria \num{12,4\%} do desequilíbrio na distribuição de
    médicos --- a \num{US\$ 15,7 milhões por ponto percentual}. Reservar vagas
    nas faculdades de medicina para quem nasceu nessas regiões corrigiria
    \num{63,8\%}, por \num{US\$ 2,2 a 5,1 milhões} o ponto.
  \end{block}

  \fonte{\textbf{Fontes:} Scott et al. (2013), \textit{Social Science \&
  Medicine}; Costa, Nunes \& Sanches (2024), \textit{Review of Economics and
  Statistics}.}
\end{frame}

\begin{frame}{Mas é caro, e não segura: a evidência contra}
  \vspace{0.4em}
  \begin{block}{E o médico vai embora quando a obrigação acaba}
    \small
    Nos \textbf{Estados Unidos}, oito anos depois, \num{12\%} dos médicos que
    foram para clínicas rurais com bolsa e \textbf{obrigação de permanência}
    ainda estavam lá --- contra \num{39\%} dos que foram \textbf{sem obrigação
    nenhuma}.
  \end{block}

  \vspace{1.2em}
  \centering
  \begin{beamercolorbox}[wd=0.92\textwidth,sep=0.85em,rounded=true,shadow=true]{block title}
    \centering Dinheiro move alocação, mas é caro, não move todo mundo, e não
    garante que quem foi fique.\\[0.35em]
    A evidência \textbf{não decide} se um degrau de R\$ 5 mil basta.
  \end{beamercolorbox}

  \fonte{\textbf{Fonte:} Pathman, Konrad \& Ricketts (1992), \textit{JAMA}.}
\end{frame}

% ===========================================================================
% 10 — No primeiro ciclo, a bolsa maior não ordenou...    (2 frames)
% ===========================================================================
\rastreio{1. Motivação \textperiodcentered\ Efeito incerto \textperiodcentered\ 3 de 3}

\begin{frame}{No primeiro ciclo, a bolsa maior não ordenou o preenchimento}
  \vspace{0.2em}
  {\small Das \num{1.295 vagas} da primeira chamada, \num{30\%} tiveram alguém
  confirmado ou homologado.}

  \vspace{0.3em}
  \begin{center}
    \includegraphics[width=\textwidth,height=0.66\textheight,keepaspectratio]%
      {preenchimento_ciclo1}
  \end{center}

  \fonte{\textbf{Fontes:} quadro de vagas e resultados do ciclo 1, chamada 1
  (Ministério da Saúde, 2025); estratos pela REGIC 2018 e pela composição de
  regiões metropolitanas e RIDEs de 2022 (IBGE).}
\end{frame}

\begin{frame}{No primeiro ciclo, a bolsa maior não ordenou o preenchimento}
  \vspace{0.4em}
  {\small A bolsa maior não veio acompanhada de mais preenchimento --- e o
  território, sim, ordenou o resultado:}

  \vspace{0.9em}
  {\small
  \begin{itemize}
    \setlength{\itemsep}{0.9em}
    \item Por \textbf{faixa de bolsa}: \lead{pagar o dobro não preencheu mais
      que pagar uma vez e meia.}
    \item Por \textbf{território}: o preenchimento cai do município
      metropolitano e da capital para o interior conectado a um polo, e é
      \lead{menor no interior remoto.}
  \end{itemize}}

  \vspace{1.1em}
  \centering
  \begin{beamercolorbox}[wd=0.92\textwidth,sep=0.8em,rounded=true]{block body}
    \small\centering Isso é \textbf{descrição, não efeito}: as faixas diferem
    em muito mais do que no valor da bolsa. Mas é o suficiente para colocar a
    pergunta.
  \end{beamercolorbox}
\end{frame}

% ===========================================================================
% 11 — Pergunta                                            (1 frame)
% ===========================================================================
\rastreio{2. Pergunta}

\begin{frame}{Pergunta}
  \vspace{0.8em}
  \centering
  \begin{beamercolorbox}[wd=0.95\textwidth,sep=1.0em,rounded=true,shadow=true]{block title}
    \centering\large Maiores bolsas do PMM-E para municípios mais vulneráveis
    compensam suas desvantagens territoriais na atração de médicos
    especialistas?
  \end{beamercolorbox}

  \vspace{1.3em}
  \raggedright
  {\small Dois objetos, um contra o outro:}

  \vspace{0.6em}
  \begin{columns}[T]
    \begin{column}{0.47\textwidth}
      \begin{beamercolorbox}[wd=\textwidth,sep=0.75em,rounded=true]{block body}
        \begin{minipage}[t][3.0em][c]{\dimexpr\textwidth-1.5em\relax}
          \small\textbf{O preço} que a política colocou sobre a
          vulnerabilidade --- o degrau de \textbf{R\$ 5 mil} entre faixas
        \end{minipage}
      \end{beamercolorbox}
    \end{column}
    \begin{column}{0.47\textwidth}
      \begin{beamercolorbox}[wd=\textwidth,sep=0.75em,rounded=true]{block body}
        \begin{minipage}[t][3.0em][c]{\dimexpr\textwidth-1.5em\relax}
          \small\textbf{A desvantagem} que esse preço pretende compensar
        \end{minipage}
      \end{beamercolorbox}
    \end{column}
  \end{columns}

  \vspace{1.0em}
  {\small A margem que observamos é o \lead{preenchimento da vaga}: se, ao
  final da chamada, apareceu alguém disposto a ocupá-la.}
\end{frame}

% ===========================================================================
% 12 — De onde vem o modelo                                (2 frames)
% ===========================================================================
\rastreio{3. Literatura teórica}

\begin{frame}{De onde vem o modelo}
  \vspace{0.3em}
  {\small O modelo junta \textbf{três tradições}. A primeira dá a
  \lead{estrutura da decisão}: o médico compara lugares pelo rendimento real
  que cada um oferece, descontado do que custa viver ali.}

  \vspace{1.0em}
  \begin{beamercolorbox}[wd=\textwidth,sep=0.85em,rounded=true]{block body}
    \small
    \textbf{Moehling, Niemesh, Thomasson \& Treber (2020)}, eq. 1, p. 184
    \\[0.45em]
    O médico escolhe a localidade que maximiza o valor presente do rendimento
    real, líquido do custo não pecuniário de viver ali.
    \\[0.3em]
    \[
      \arg\max_{i \in I}\ \sum_t \delta^t
      \left[ \dfrac{\mathbb{E}\!\left(w_{it}^{(s)}\right)}{p_{it}}
      - c_{it}^{(s)} \right]
    \]
  \end{beamercolorbox}
\end{frame}

\begin{frame}{De onde vem o modelo}
  \vspace{0.1em}
  {\small As outras duas \textbf{abrem o custo}: uma diz o que há dentro do
  \lead{custo geográfico}, a outra o que há dentro do \lead{custo de
  trabalhar}.}

  \vspace{0.55em}
  \begin{beamercolorbox}[wd=\textwidth,sep=0.6em,rounded=true]{block body}
    \small
    \textbf{Redding \& Rossi-Hansberg (2017)}, eq. 24, p. 28 \quad
    {\color{cinzaFonte}custo geográfico}
    \\[0.3em]
    A utilidade de trabalhar em um lugar depende do salário, das amenidades e
    do custo de moradia locais.
    \\[-0.2em]
    \[ u_{nio} = \dfrac{z_{nio}\, B_n\, w_i}{\kappa_{ni}\, Q_n^{1-\beta}} \]
  \end{beamercolorbox}

  \vspace{0.45em}
  \begin{beamercolorbox}[wd=\textwidth,sep=0.6em,rounded=true]{block body}
    \small
    \textbf{Choné \& Ma (2011)}, eq. 1, p. 232 \quad
    {\color{cinzaFonte}custo de trabalhar}
    \\[0.3em]
    A utilidade do médico soma a renda, subtrai o custo de atender e soma o
    benefício ao paciente, ponderado pelo altruísmo. Equipe e capital instalado
    entram no custo de atender.
    \\[-0.2em]
    \[ U = R - C(q; L, K) + \alpha B(q) \]
  \end{beamercolorbox}

  \vspace{0.35em}
  {\small Nenhuma das três trata de um componente da remuneração fixado por
  \lead{regra pública sobre um índice territorial}. É isso que a adaptação ao
  PMM-E acrescenta.}
\end{frame}

% ===========================================================================
% 13 — Como o médico escolhe onde trabalhar   (1 frame, 2 overlays)
% A equacao fica fixa; primeiro a condicao de aceitacao, depois o glossario.
% ===========================================================================
\rastreio{4. Modelo microeconômico \textperiodcentered\ 1 de 4}

\begin{frame}{Como o médico escolhe onde trabalhar}
  \vspace{0.3em}
  {\small Adaptando Moehling et al. (2020), o médico $i$ escolhe o município
  $m$ que maximiza o \lead{valor presente líquido da carreira}:}

  \vspace{-0.2em}
  \[
    V_{im} = \sum_t \delta^t
    \left[ \frac{\mathbb{E}(w_{imt} \mid B_m)}{p_{mt}} - c_{im} \right]
    + \varepsilon_{im},
    \qquad
    m_i^{\ast} \in \arg\max_{m \in \mathcal{M} \cup \{0\}} V_{im}
  \]

  \only<1>{%
    \vspace{1.2em}
    \begin{beamercolorbox}[wd=\textwidth,sep=0.85em,rounded=true]{block body}
      \small
      A alternativa $m = 0$ é \textbf{ficar fora do programa}. A vaga em $m$ é
      aceita quando $V_{im} \geq V_{i0}$ e $m$ é a melhor opção entre as
      disponíveis.
    \end{beamercolorbox}
  }%
  \only<2>{%
    \vspace{0.3em}
    {\footnotesize\renewcommand{\arraystretch}{1.12}
    \begin{tabularx}{\textwidth}{@{}>{\centering\arraybackslash}p{0.22\textwidth} L@{}}
      \toprule
      \bfseries Termo & \bfseries Leitura \\
      \midrule
      $\mathbb{E}(w_{imt} \mid B_m)$ & remuneração esperada no município, dada
        a bolsa $B_m$ fixada pela regra \\
      $p_{mt}$ & custo de vida local: o que importa é a remuneração
        \textbf{real} \\
      $c_{im}$ & tudo o que torna estar naquele município custoso e não é pago
        em dinheiro \\
      $\delta^t$ & a decisão é de carreira, não de um mês: o médico pesa o que
        o lugar oferece ao longo do tempo \\
      $\varepsilon_{im}$ & gostos individuais que não observamos \\
      \bottomrule
    \end{tabularx}}
  }%
\end{frame}

% ===========================================================================
% 14 — O que a bolsa paga — e o que não paga               (2 frames)
% ===========================================================================
\rastreio{4. Modelo microeconômico \textperiodcentered\ 2 de 4}

\begin{frame}{O que a bolsa paga --- e o que não paga}
  \vspace{0.4em}
  {\small A remuneração tem duas partes: a \textbf{bolsa}, que a regra fixa, e
  o que o médico obtém no \lead{mercado local} fora das 20 horas do programa:}

  \[ \mathbb{E}(w_{imt} \mid B_m) = B_m + w^{\text{priv}}_m \]

  \vspace{0.5em}
  {\small
  \begin{tabularx}{\textwidth}{@{}>{\raggedright\arraybackslash}p{0.26\textwidth} L >{\centering\arraybackslash}p{0.29\textwidth}@{}}
    \toprule
    \bfseries Onde & \bfseries Mercado privado & \bfseries Remuneração \\
    \midrule
    Capital ou região metropolitana & consultório, planos de saúde, hospitais
      privados & \makecell{$w = B + w^{\text{priv}}$,\\ com $w^{\text{priv}}$ alto} \\
    \addlinespace[0.2em]
    Interior isolado & sem demanda privada que sustente a especialidade &
      $w \to B$ \\
    \bottomrule
  \end{tabularx}}
\end{frame}

\begin{frame}{O que a bolsa paga --- e o que não paga}
  \vspace{0.5em}
  \lead{Duas consequências}

  \vspace{0.8em}
  \begin{block}{Primeiro}
    \small
    A bolsa maior do interior pode significar \textbf{remuneração total
    menor}: R\$ 20 mil sem complemento contra R\$ 10 mil mais o consultório da
    capital. Para compensar, $B$ precisa cobrir também a renda privada que o
    médico deixa de ganhar.
  \end{block}

  \vspace{0.6em}
  \begin{block}{Segundo}
    \small
    Onde $w^{\text{priv}} \approx 0$ a bolsa é \textbf{toda} a remuneração ---
    e é exatamente ali que a política mais aposta nela.
  \end{block}

  \vspace{1.0em}
  {\small O deflator $p_{mt}$ trabalha no sentido oposto: o custo de vida é
  \textbf{menor} no interior, o que valoriza a mesma bolsa em termos reais.}
\end{frame}

% ===========================================================================
% 15 — O custo de estar ali                                (3 frames)
% ===========================================================================
\rastreio{4. Modelo microeconômico \textperiodcentered\ 3 de 4}

\begin{frame}{O custo de estar ali}
  \vspace{0.3em}
  \[
    c_{im} =
    \underbrace{\phi(\text{dist}_{im}) - \gamma A_m + \theta_i^{\text{rural}}}%
      _{\text{geográfico}}
    + \underbrace{C(q; L_m, K_m) - \alpha_i B(q; L_m, K_m)}%
      _{\text{laboral líquido}}
  \]

  \vspace{0.3em}
  {\small
  \begin{tabularx}{\textwidth}{@{}>{\centering\arraybackslash}p{0.20\textwidth} >{\centering\arraybackslash}p{0.20\textwidth} L@{}}
    \toprule
    \bfseries Componente & \bfseries Sinal & \bfseries Significado \\
    \midrule
    $\phi(\text{dist})$ & $\phi' > 0$ & afastar-se da família custa, e custa
      mais a cada quilômetro \\
    $-\gamma A_m$ & $< 0$ & amenidade urbana --- saneamento, segurança,
      escola --- compensa \\
    $\theta_i^{\text{rural}}$ & $\gtrless 0$ & gosto pessoal por cidade pequena
      ou grande \\
    $C(q)$ & $C' > 0,\ C'' > 0$ & atender cansa, e cansa de forma crescente \\
    $\alpha_i B(q)$ & $B' > 0,\ B'' < 0$ & curar dá satisfação, ponderada pelo
      altruísmo $\alpha_i$ \\
    \bottomrule
  \end{tabularx}}
\end{frame}

\begin{frame}{O custo de estar ali}
  \vspace{0.4em}
  {\small O bloco \textbf{geográfico} vem de Redding \& Rossi-Hansberg; o
  \textbf{laboral}, de Choné \& Ma.}

  \vspace{0.8em}
  {\small Equipe ($L$) e capital instalado ($K$) atuam \lead{duas vezes} no
  bloco laboral:}

  \vspace{0.5em}
  \begin{columns}[T]
    \begin{column}{0.47\textwidth}
      \begin{beamercolorbox}[wd=\textwidth,sep=0.75em,rounded=true]{block body}
        \begin{minipage}[t][4.4em][c]{\dimexpr\textwidth-1.6em\relax}
          \small\centering reduzem o cansaço\\[0.4em]
          $\dfrac{\partial C}{\partial K} < 0$
        \end{minipage}
      \end{beamercolorbox}
    \end{column}
    \begin{column}{0.47\textwidth}
      \begin{beamercolorbox}[wd=\textwidth,sep=0.75em,rounded=true]{block body}
        \begin{minipage}[t][4.4em][c]{\dimexpr\textwidth-1.6em\relax}
          \small\centering ampliam o benefício que o atendimento
          produz\\[0.4em]
          $\dfrac{\partial B}{\partial K} > 0$
        \end{minipage}
      \end{beamercolorbox}
    \end{column}
  \end{columns}

  \vspace{0.7em}
  {\small\centering Por dois caminhos, $\partial c / \partial K < 0$. A
  segunda é \textbf{extensão deste projeto} a Choné \& Ma.\par}

  \vspace{0.9em}
  \centering
  \begin{beamercolorbox}[wd=0.95\textwidth,sep=0.8em,rounded=true,shadow=true]{block title}
    \centering Como no município vulnerável $L$ e $K$ são baixos, o mesmo lugar
    que paga mais é o que impõe maior custo de trabalho.
  \end{beamercolorbox}
\end{frame}

\begin{frame}{O custo de estar ali}
  \vspace{0.2em}
  \begin{center}
    \includegraphics[width=\textwidth,height=0.74\textheight,keepaspectratio]%
      {curva_custo_laboral_burnout}
  \end{center}

  \fonte{\textbf{Fonte:} ilustração conceitual do modelo --- custo laboral
  líquido em função do volume de atendimentos.}
\end{frame}

% ===========================================================================
% 16 — O IVS organiza o custo                              (2 frames)
% ===========================================================================
\rastreio{4. Modelo microeconômico \textperiodcentered\ 4 de 4}

\begin{frame}{O IVS organiza o custo}
  \vspace{0.1em}
  {\small Distância da família, aluguel e esforço clínico \textbf{não são
  observados}. O que se observa, para todo município, é o \lead{IVS}.
  Escrevemos então o custo como função do índice mais um desvio individual:}

  \vspace{-0.4em}
  \[ c_{im} = c_0(IVS_m) + \eta_i \]
  \vspace{-1.0em}

  {\small Isso não é atalho: \lead{cada dimensão do IVS corresponde a um bloco
  do custo.}}

  \vspace{0.4em}
  {\footnotesize\renewcommand{\arraystretch}{1.15}
  \begin{tabularx}{\textwidth}{@{}>{\raggedright\arraybackslash}p{0.17\textwidth} >{\raggedright\arraybackslash}p{0.23\textwidth} L >{\centering\arraybackslash}p{0.13\textwidth}@{}}
    \toprule
    \bfseries Dimensão do IVS & \bfseries Indicadores & \bfseries Bloco do
      custo & \bfseries Efeito sobre $c$ \\
    \midrule
    Infraestrutura urbana & saneamento, lixo, tempo de deslocamento &
      amenidades $A_m$ & $\uparrow$ \\
    Renda e trabalho & pobreza, desemprego, informalidade & mercado privado
      ausente, $w \to B$ & $\uparrow$ \\
    Capital humano & mortalidade infantil, analfabetismo, mães adolescentes &
      gravidade do caso, $B'(q)\uparrow$; e escassez de equipe e capital,
      $K\downarrow$ & \textbf{ambíguo} \\
    \bottomrule
  \end{tabularx}}
\end{frame}

\begin{frame}{O IVS organiza o custo}
  \vspace{0.4em}
  {\small \lead{A terceira linha é o que impede assumir que o custo cresce com
  o índice.}}

  \vspace{0.8em}
  \begin{columns}[T]
    \begin{column}{0.47\textwidth}
      \begin{beamercolorbox}[wd=\textwidth,sep=0.8em,rounded=true]{block body}
        \begin{minipage}[t][3.4em][c]{\dimexpr\textwidth-1.7em\relax}
          \hyphenpenalty=10000\exhyphenpenalty=10000
          \small Carência sanitária \textbf{eleva o benefício} de atender ---
          o que \textbf{reduz} o custo para um médico altruísta.
        \end{minipage}
      \end{beamercolorbox}
    \end{column}
    \begin{column}{0.47\textwidth}
      \begin{beamercolorbox}[wd=\textwidth,sep=0.8em,rounded=true]{block body}
        \begin{minipage}[t][3.4em][c]{\dimexpr\textwidth-1.7em\relax}
          \hyphenpenalty=10000\exhyphenpenalty=10000
          \small Ao mesmo tempo \textbf{sinaliza falta de insumos} --- o que
          \textbf{eleva} o cansaço.
        \end{minipage}
      \end{beamercolorbox}
    \end{column}
  \end{columns}

  \vspace{0.9em}
  \centering
  {\large $c_0'(IVS) \gtrless 0$ \quad {\normalsize o sinal é questão
  empírica}}

  \vspace{1.1em}
  \begin{beamercolorbox}[wd=0.93\textwidth,sep=0.85em,rounded=true,shadow=true]{block title}
    \centering É por isso que o objeto do trabalho é o \textbf{degrau} da bolsa
    entre faixas, e não a inclinação do índice.
  \end{beamercolorbox}
\end{frame}

% ===========================================================================
% 17 — Duas hipóteses                                      (2 frames)
% ===========================================================================
\rastreio{5. Hipóteses}

\begin{frame}{Duas hipóteses}
  \vspace{0.5em}
  \lead{Passo 1 --- a condição de aceitação.}
  {\small O médico $i$ aceita a vaga em $m$ quando o que ela vale supera sua
  melhor alternativa $\bar{v}_i$:}

  \vspace{0.2em}
  \[
    \frac{B_m + w^{\text{priv}}_m}{p_m} - c_0(IVS_m) \geq \bar{v}_i
  \]

  \vspace{1.0em}
  \lead{Passo 2 --- do médico à vaga.}

  \vspace{0.4em}
  \begin{beamercolorbox}[wd=\textwidth,sep=0.85em,rounded=true]{block body}
    \small A vaga é preenchida se existir \textbf{ao menos um} candidato para
    quem a desigualdade vale. Tudo o que aumenta o lado esquerdo aumenta essa
    probabilidade.
  \end{beamercolorbox}
\end{frame}

\begin{frame}{Duas hipóteses}
  \vspace{0.2em}
  \lead{Passo 3 --- as derivadas.}

  \vspace{0.4em}
  {\footnotesize\renewcommand{\arraystretch}{1.1}
  \begin{tabularx}{\textwidth}{@{}>{\centering\arraybackslash}p{0.06\textwidth} L >{\centering\arraybackslash}p{0.33\textwidth}@{}}
    \toprule
    & \bfseries Hipótese & \bfseries Derivada \\
    \midrule
    \textbf{H1} & Maior \textbf{remuneração real} aumenta a probabilidade de
      preenchimento da vaga
      & $\dfrac{\partial \Pr(\text{preench.}_m)}{\partial (B_m / p_m)} > 0$ \\
    \addlinespace[0.5em]
    \textbf{H2} & Maior \textbf{custo locacional} reduz a probabilidade de
      preenchimento da vaga
      & $\dfrac{\partial \Pr(\text{preench.}_m)}{\partial c_m} < 0$ \\
    \bottomrule
  \end{tabularx}}

  \vspace{1.0em}
  \lead{Passo 4 --- as duas juntas.}
  {\small Na fronteira entre duas faixas, o preenchimento do lado mais
  vulnerável exige que o \textbf{degrau monetário} supere o \textbf{degrau de
  custo}:}

  \vspace{-0.2em}
  \[ \frac{\Delta B_m}{p_m} > \Delta c_0,
     \qquad \Delta B_m = \text{R\$ }5.000 \]
  \vspace{-0.8em}

  \centering{\small A pergunta da apresentação é se essa desigualdade vale.\par}
\end{frame}

% ===========================================================================
% 18 — Viabilidade empírica                                (2 frames)
% ===========================================================================
\rastreio{6. Viabilidade empírica}

\begin{frame}{Viabilidade empírica}
  \vspace{0.2em}
  \lead{Temos como medir cada peça do modelo?}

  \vspace{0.5em}
  {\footnotesize
  \begin{tabularx}{\textwidth}{@{}>{\raggedright\arraybackslash}p{0.20\textwidth} L >{\raggedright\arraybackslash}p{0.28\textwidth}@{}}
    \toprule
    \bfseries Peça do modelo & \bfseries O que observamos & \bfseries Fonte \\
    \midrule
    Preenchimento da vaga & se a vaga teve alguém confirmado ou homologado &
      quadro de vagas e resultados do ciclo 1: 1.295 vagas em 368 municípios \\
    Bolsa $B_m$ & o valor anunciado na vaga: R\$ 10, 15 ou 20 mil & edital \\
    Custo $c_m$ & o IVS e suas três dimensões; se o município é capital,
      metropolitano, polo do interior ou interior remoto; quantos especialistas
      e que estrutura já havia & Ipea; REGIC 2018 e RMs 2022 (IBGE); CNES
      mensal, jun/2024 a jul/2026 \\
    Custo de vida $p_m$ & diferenças entre estados & IBGE \\
    Mercado local $w^{\text{priv}}_m$ & --- & \textbf{não observado} \\
    Distância da família & --- & \textbf{não observado} \\
    \bottomrule
  \end{tabularx}}
\end{frame}

\begin{frame}{Viabilidade empírica}
  \vspace{0.4em}
  {\small \lead{Sim, para o essencial.} Duas peças ficam de fora --- a renda no
  mercado privado local e a distância da família --- e entram no modelo como
  parte do custo que o IVS e a tipologia territorial resumem.}

  \vspace{1.0em}
  \begin{block}{A dificuldade}
    \small
    A bolsa e a vulnerabilidade andam juntas porque a regra fixa uma pela
    outra: \num{não existe município com bolsa alta e vulnerabilidade baixa}.

    \vspace{0.45em}
    Separar H1 de H2 exige comparar municípios parecidos que caíram em faixas
    diferentes --- e a faixa publicada não coincide com a que sai dos cortes do
    Ipea em \textbf{177 dos 368 municípios} do ciclo 1.
  \end{block}

  \vspace{0.9em}
  \centering
  \begin{beamercolorbox}[wd=0.93\textwidth,sep=0.85em,rounded=true,shadow=true]{block title}
    \centering Recuperar a regra exata que o Ministério aplicou é o primeiro
    passo do trabalho empírico.
  \end{beamercolorbox}
\end{frame}

\end{document}
